% ==========================================================================
% Paper 2 — Section 5: The Geometric Boundary of Routing. Hyperplane
% Bisection in d=2 (TAN-II Sprints 4.2-A / 4.2-B, frozen).
% ==========================================================================
\section{The Geometric Boundary of Routing: Hyperplane Bisection in
\texorpdfstring{$d=2$}{d=2}}
\label{sec:routing}

Section 4 closed the dynamical avenue: opponent competition reshapes
tuning and gain but cannot reorder candidates, because ordering is decided
by the kernel before any downstream dynamics. The boundary is therefore
representational, and this section crosses it with the minimal
representational change---operator O2 of §2.4, the vector query--key
interaction. The section proceeds in four parts: the exact theorem that
explains why the scalar lineage was locked (the order-conservation
theorem); the pure geometry of hyperplane bisection and its exact flip
law $\mathrm{FlipRate}_{\mathrm{CF}}(\theta)=\theta/\pi$; the
ordering-capacity ladder that generalizes the flip law beyond two
candidates; and the realization of all of this inside the running neuron
(Sprint 4.2-B), including the shortcut audits that certify the flips are
not artifacts.

\subsection{The order-conservation theorem}
\label{sec:orderconservation}

\begin{theorem}[Order conservation of the separable scalar kernel]
\label{thm:order}
Let $d=1$, $K_A\neq K_B$, and let the query domain $D\subseteq\mathbb{R}$
intersect at most one sign class of $\mathbb{R}\setminus\{0\}$. Then the
winner of the dot-product kernel $E_i(Q)=QK_i$ is constant on $D$:
$w(Q)=\operatorname{argmax}_i K_i$ for $D\subseteq(0,\infty)$ and
$w(Q)=\operatorname{argmin}_i K_i$ for $D\subseteq(-\infty,0)$. In
particular, for the frozen TAN kernel $E_{t,i}=\beta W_qW_kS_tx_i$ the
query $q_t=\beta W_qW_kS_t$ lies in the closed cone $\mathcal{Q}_+=
[0,\infty)$, so $\operatorname{argmax}_i E_{t,i}=\operatorname{argmax}_i
x_i$ for every active query, and the counterfactual flip rate is exactly
zero.
\end{theorem}

\begin{proof}
$QK_A>QK_B\iff (Q>0 \wedge K_A>K_B)\vee(Q<0\wedge K_A<K_B)$, so the winner
is constant on each sign class and changes only when the query crosses
zero. The \TAN{} query never does: $S_t=[x_t-\mu_t-\varepsilon]_+\ge0$ and
$\beta,W_q,W_k>0$.
\end{proof}

Theorem~\ref{thm:order} is the paper's sharpest statement of the wall that
Sprint 4.2-A had to audit, and it is deliberately \emph{domain-conditional}:
a signed scalar query ($D=\{+1,-1\}$) \emph{can} flip the winner, but only
by the degenerate global reversal $w(-1)=\operatorname{argmin}_iK_i$, a
one-bit, order-locked capacity (flip rates confined to $\{0,1\}$). The
correct universal claim is therefore never ``one-dimensional attention
cannot route'' but: \emph{within the positive-query scalar family inherited
from TAN-I, candidate ordering is query-invariant}. This scoping is frozen
as Amendment AM1 of the programme and is used verbatim here.

\subsection{Hyperplane bisection and the flip law}
\label{sec:bisection}

Fix two candidates $K_A,K_B\in\mathbb{R}^d$ ($d\ge2$) with
$\Delta K=K_A-K_B\neq0$, and sample the query isotropically,
$Q\sim\mathrm{Unif}(S^{d-1})$. The decision boundary
$\mathcal{H}_0=\{Q:Q^{\top}\Delta K=0\}$ is a hyperplane through the
origin, i.e.\ a great $(d-2)$-sphere on $S^{d-1}$, so it bisects the query
sphere into two open hemispheres of equal measure---\emph{regardless of the
candidate norms}, because the boundary is homogeneous.

\begin{proposition}[Equal bisection]
\label{prop:bisection}
For any $K_A\neq K_B$ and isotropic $Q$ on $S^{d-1}$,
$\Pr_Q[w(Q)=A]=\Pr_Q[w(Q)=B]=\tfrac12$.
\end{proposition}

\begin{proof}
$\Pr[Q^{\top}\Delta K>0]=\Pr[Q\cdot\hat u>0]$ with
$\hat u=\Delta K/\|\Delta K\|$; the projection of a uniform unit vector
onto a fixed axis is symmetric about zero.
\end{proof}

\begin{corollary}[Dimension-independent winner variation]
\label{cor:fd}
For i.i.d.\ queries $Q_1,Q_2$,
$F_d=\Pr[w(Q_1)\neq w(Q_2)]=2\cdot\tfrac12\cdot\tfrac12=\tfrac12$ for
every $d\ge2$; the same holds for the signed scalar ($d=1$, uniform on
$\{\pm1\}$), while the positive scalar lineage gives $F_d=0$. The Monte
Carlo estimates freeze at $0.4996/0.4997/0.5033/0.4998$ for
$d=2,3,4,8$: \emph{higher dimension does not mean stronger routing}, and
$F_d$ is retained in this paper only as a negative control, never as a
transition statistic.
\end{corollary}

The counterfactual quantity that \emph{does} carry the geometry is the flip
rate of a fixed query pair over randomized histories. Let the history be
isotropic, $\Delta K\sim\mathrm{Unif}(S^{d-1})$, and fix two queries
$Q_A,Q_B\in S^{d-1}$ with angular separation
$\theta=\arccos(Q_A^{\top}Q_B)\in[0,\pi]$.

\begin{proposition}[The flip law]
\label{prop:fliplaw}
$\mathrm{FlipRate}_{\mathrm{CF}}(\theta)
=\Pr_{\Delta K}\bigl[w(Q_A)\neq w(Q_B)\bigr]=\theta/\pi$.
\end{proposition}

\begin{proof}
Rotate so that $Q_A=e_1$, $Q_B=\cos\theta\,e_1+\sin\theta\,e_2$. A flip
occurs iff $(\Delta K_1)\bigl(\cos\theta\,\Delta K_1+\sin\theta\,\Delta
K_2\bigr)<0$. Writing $\Delta K_1=r\cos\phi$, $\Delta K_2=r\sin\phi$, the
marginal law of $(\Delta K_1,\Delta K_2)$ is rotationally symmetric in the
$(1,2)$-plane, so $\phi$ is uniform on $[0,2\pi)$, and the event is
$\cos\phi\cos(\phi-\theta)<0$, the union of two arcs of angular width
$\theta$: total measure $2\theta$ out of $2\pi$.
\end{proof}

The frozen verification is exact at both levels. Analytically, a
root-decomposition of the indicator reproduces $\theta/\pi$ on the full
11-point angle grid (endpoints $0\to0$, $\pi\to1$) at machine precision,
independently of the closed form (connecting to the classical random
hyperplane flip law of \citet{charikar2002}); and the Monte Carlo
audit---$d\in\{2,3,4,8\}\times11$ angles $\times3$ seeds, $132/132$
conditions, maximum deviation across 3-seed averages $0.004667$ (maximum
single-seed absolute deviation $0.009267$; grid parameter recorded as
$\theta/\pi$)---agrees with the curve point by point. The
degenerate endpoints of the scalar lineage sit exactly where the theorem
predicts: the signed scalar realizes only $\{0,1\}$ (queries equal or
opposite), the positive scalar only $\{0\}$. The transition recorded in
this section is therefore not ``$d=1\to d=2$ increases flips'' but the
topological change $S^{0}\to S^{1}$: the query-direction manifold becomes
connected, and the flip rate becomes a continuous function
$\theta/\pi\in[0,1]$ of a continuous query angle.

\subsection{The ordering-capacity ladder}
\label{sec:capacity}

The pairwise flip law extends to a capacity statement over $M$ candidates.
For $M$ keys in general position, the achievable orderings of
$(Q^{\top}K_1,\ldots,Q^{\top}K_M)$ as $Q$ ranges over the query domain are
the cells of the central arrangement of the $C(M,2)$ difference
hyperplanes. Frozen exact computations (arrangement enumeration, linear
programming existence, and a $200{,}000$-angle dense sampling cross-check)
establish the ladder:

\begin{center}
\begin{tabular}{@{}lcccc@{}}
\toprule
$M$ & positive scalar & signed scalar & $d=2$ & $d=3,8$ \\
\midrule
2 & 1 & 2 & 2 & 2 \\
3 & 1 & 2 & 6 & 6 \\
4 & 1 & 2 & 12 & 24 \\
\bottomrule
\end{tabular}
\end{center}

\noindent i.e.\ $1$ / $2$ / $M(M-1)$ / $M!$ (the last for $M\le d+1$ in
general position). Two facts anchor the paper's reading. First, the signed
scalar is stuck at two orderings for \emph{every} $M$: its ``routing'' is
the global reversal of a fixed magnitude order, never a content address.
Second, $d=2$ already realizes \emph{all} $3!=6$ orderings of three
candidates, so $d=2$ is the minimal dimension for a connected,
rotationally symmetric continuous query-direction manifold---the statement
this paper is willing to make, with the qualifier that it concerns the
tested positive-scalar lineage and the tested query--key family. Higher
ambient dimension adds nothing to this pairwise or three-candidate
capacity.

\subsection{Realization in the running neuron (Sprint 4.2-B)}
\label{sec:realization}

The geometry is only half of the rung; the frozen programme demands that
the running model actually flip, under the iron gate that \emph{only the
representation changes} (window, surprise gate, membrane, readout all
frozen). Sprint 4.2-B does exactly this with operator O2
($\hat\varphi_2$ keys, rotating query $c\,S_t\hat u(\omega S_t)$,
$\omega=0.4$), and its deterministic canonical audit closed every
pre-registered prediction:

\begin{itemize}
\item \emph{Counterfactual winner flip.} Fixed history
$[1,0,2,0]$, only the query changes $3\to4$: the attention winner flips
$A\to B$ with energy margins $+0.2587427\to-0.1559101$; the crossing
$Q^*=3.7071$ matches the closed form and the winner-vs-query scan
(A for $Q\le3.70$, B for $Q\ge3.75$).
\item \emph{No amplitude-code leakage.} The probe slot's energy
($3.1078$ at $Q=3$; $4.9890$ at $Q=4$) stays \emph{below} the winning
event at both queries, in direct contrast to the scalar lineage, whose
probe dominates the window ($10.8$/$20.8$).
\item \emph{The JSD conjunction.} In a 2-candidate geometric toy example,
the scalar trap exhibits JSD$=0.0751$ with zero flip (query changes
sharpness, not order), while the vector model exhibits JSD$=0.0647$ with
a winner flip. In the full 5-slot window ($[1,0,2,0,Q]$), the canonical
distribution shifts yield full-window JSDs of $0.009462$ (M0) and $0.008727$ (M2).
\item \emph{Normalization is a mechanistic enabler.} With raw magnitude
keys the kernel reads
$(x_A-x_B)[\cos\theta+(x_A+x_B)\sin\theta]$, whose sign is fixed in the
query range: the unnormalized control has flip rate exactly $0$ (frozen
Amendment AM2), while the normalized model reaches $0.4425$ on the
randomized ensemble against the $0.4445$ quadrature benchmark
(3 seeds $\times10^4$, all within criterion).
\item \emph{Structure, not luck.} The sweep over
$\omega\in\{0.2,\allowbreak0.4,\allowbreak\pi/2,\allowbreak\pi\}$ yields
winners $(A,A)$, $(A,B)$, $(B,B)$, $(A,B)$: the flip exists exactly when
the query direction sweeps across the decision boundary, with the
closed-form crossings in place.
\item \emph{Dimension controls.} In this mechanism family, increasing the
ambient embedding dimension beyond the minimal 2-D query subspace did not
increase routing capacity and, under the tested finite
candidate/query configuration, \emph{reduced} the observed flip rate
(ensemble benchmarks $0.0549$/$0.0284$/$0.0172$ for $d=3,4,8$; canonical
flips absent). The cause is a representation--candidate geometry
mismatch---the key-difference vector leaves the 2-D query subspace and the
attention mass is absorbed by the silent slots---not a ``curse of
dimensionality''.
\end{itemize}

\begin{figure}[t]
\centering
\includegraphics[width=0.66\textwidth]{fig3_phase}\hfill
\includegraphics[width=0.32\textwidth]{fig3d_realization}
\caption{The geometric boundary of routing. (a)--(c) The pure geometry
(Sprint 4.2-A): the signed scalar ray realizes only the degenerate flip
rates $\{0,1\}$; in $d=2$ the decision hyperplane bisects the query sphere;
and the Monte Carlo flip rates sit on the analytic curve
$\mathrm{FlipRate}_{\mathrm{CF}}(\theta)=\theta/\pi$ (132/132 conditions,
$d\in\{2,3,4,8\}$). (d) The realization in the running neuron
(Sprint 4.2-B): on the fixed history $[1,0,2,0]$, changing only the query
flips the attention winner at the closed-form crossing $Q^*=3.7071$, with
margins $+0.2587\to-0.1559$.}
\label{fig:phase}
\end{figure}

\subsection{Sufficiency of 1D Metric Compatibility and the Multi-Candidate Address Identifiability Boundary}
\label{sec:metric_sufficiency}

The analysis in \S\ref{sec:orderconservation}--\ref{sec:bisection} demonstrated that within the scalar dot-product family inherited from TAN-I, candidate ordering is invariant under positive queries, and that two-dimensional rotation ($S^1$) unlocks counterfactual winner flips. This raises a fundamental theoretical and architectural question: \emph{is multidimensional representation strictly necessary for single-neuron content addressing, or was the scalar failure an artifact of separable dot-product multiplication ($E(q, k) = q \cdot k$)?}

To resolve this question definitively, we established two formal propositions and executed a comprehensive multi-candidate address identifiability audit (60,000 synthetic histories, 420,000 model evaluations across 7 model architectures and 6 conditions).

\begin{proposition}[1D Metric Sufficiency under Hard Selection]
\label{prop:metric_suff}
Let $\{k_1, k_2, \dots, k_N\} \subset \mathbb{R}$ be a set of $N \ge 2$ distinct scalar keys. Define the scalar distance compatibility kernel by $E(q, k) = -(q - k)^2$. Then for any target key $k_t$, setting the query to $q = k_t$ uniquely selects $k_t$ under hard argmax routing:
\begin{equation}
\operatorname{argmax}_{i \in \{1,\dots,N\}} E(k_t, k_i) = t.
\end{equation}
Furthermore, the selection margin is strictly positive: $\Delta E = \min_{i \ne t} (k_i - k_t)^2 > 0$.
\end{proposition}

\begin{proposition}[Impossibility of Scalar Intermediate Addressing]
\label{prop:scalar_impossible}
Let $N \ge 3$ and let $k_{(1)} < k_{(2)} < \dots < k_{(N)}$ be ordered distinct scalar keys. For any generalized scalar multiplication kernel of the form $E(q, k) = \alpha(q) \cdot k$ where $\alpha: \mathbb{R} \to \mathbb{R}$ is an arbitrary scalar function:
\begin{enumerate}[label=(\roman*)]
\item If $\alpha(q) > 0$, $\operatorname{argmax}_i E(q, k_i) = k_{(N)}$;
\item If $\alpha(q) < 0$, $\operatorname{argmax}_i E(q, k_i) = k_{(1)}$;
\item If $\alpha(q) = 0$, $E(q, k_i) \equiv 0$ for all $i$ (tie across all candidates).
\end{enumerate}
In all cases, any intermediate key $k_{(j)}$ with $1 < j < N$ cannot be selected as the unique maximum for any query $q \in \mathbb{R}$.
\end{proposition}

\begin{proof}
Direct from linearity: for any non-zero $\alpha(q)$, multiplication by $\alpha(q)$ either strictly preserves ($\alpha > 0$) or strictly reverses ($\alpha < 0$) the order of the keys. An intermediate element in a strictly ordered sequence can never become the maximum under order-preserving or order-reversing bijections.
\end{proof}

\begin{figure}[t]
\centering
\includegraphics[width=\textwidth]{fig7_p1a_identifiability}
\caption{Address identifiability and the multi-candidate routing boundary (420,000 model evaluations). (a) On $N=3$ candidates (Condition C2), intermediate keys cannot be addressed by scalar multiplier kernels: SignedScalar achieves exactly $0/6{,}705$ hits ($0.0\%$), whereas 1D Metric Attention achieves $100.00\%$ hit rate, identical to 2D Vector-QK on $S^1$. (b) Dual-channel both-routed accuracy across clean ($N=2, 3$), noise-perturbed ($\sigma \in \{0.01, 0.05\}$), and out-of-distribution amplitude ranges (C5, C6). (c) Soft addition readout error on C2: hard routing accuracy is $100\%$ for both 1D Metric and 2D Vector-QK, but soft readout reveals residual leakage governed by kernel curvature ($0.4247$ vs.\ $0.9463$).}
\label{fig:p1a_identifiability}
\end{figure}

The empirical consequences of Propositions~\ref{prop:metric_suff} and \ref{prop:scalar_impossible} were verified across 6 experimental conditions (Figure~\ref{fig:p1a_identifiability}):
\begin{itemize}
\item \emph{Empirical refutation of 2D necessity:} On Condition C1 ($N=2$) and Condition C2 ($N=3$), the 1D Metric Attention model achieves $100.00\%$ both-routed accuracy with zero hard addition error, exactly matching the 2D Vector-QK model on $S^1$ across all 60,000 trials. This proves that 2D orthogonal rotation is \emph{not} mathematically necessary for content addressing: a 1D scalar distance metric is sufficient.
\item \emph{Empirical confirmation of the intermediate-key barrier:} In Condition C2 ($N=3$, 10,000 trials $\times$ 2 query channels), exactly $6{,}705$ channel queries targeted the intermediate key ($k_{(1)} < k_{\text{target}} < k_{(3)}$). In exact agreement with Proposition~\ref{prop:scalar_impossible}, the SignedScalar kernel achieved \textbf{exactly $0$ hits out of $6{,}705$ queries ($0.00\%$ hit rate)}, whereas 1D Metric Attention and 2D Vector-QK achieved $6{,}705/6{,}705$ ($100.00\%$).
\item \emph{Demarcation of Historical 4.2-B:} The static control model \textsc{Hist4.2-B} achieved only $13.68\%$ both-routed accuracy on C2 (middle-key channel hit rate $36.78\%$, near the $33.33\%$ blind chance level). This demonstrates that the winner flips observed in historical Sprint 4.2-B were fixed counterfactual switches between two pre-tuned surprise amplitudes, rather than true content-addressable memory retrieval.
\item \emph{Soft readout error bounds:} Under soft softmax readout, the addition error on C2 differs between 1D Metric ($0.4247$) and 2D Vector-QK ($0.9463$). For a distance metric kernel, the soft error is strictly bounded by the minimum key separation $\delta$ and temperature $\gamma$:
\begin{equation}
|\hat{v}_{\mathrm{soft}} - v_t| \le M \frac{(N-1)e^{-\gamma\delta^2}}{1 + (N-1)e^{-\gamma\delta^2}}, \qquad M = \max_{i\ne t}|v_i - v_t|.
\end{equation}
\end{itemize}

\subsection{What this section does and does not establish}
\label{sec:scope5}

The scoped claim is: \emph{within the tested positive-scalar lineage, $d=2$
is the minimal dimension for continuous query-direction freedom under dot-product kernels, and the
running neuron realizes the resulting query-conditioned reordering
(Level-2 realized routing)}. However, moving from dot-product ranking to 1D distance metric compatibility ($E = -(q-k)^2$) restores full multi-candidate content addressability without requiring higher-dimensional rotation. It is \emph{not} claimed that $d=2$ routes in
every representation family (the $d=3,4,8$ controls show the opposite for
the tested maps), nor that any flip constitutes semantic binding---that
requires the key--value decoupling of the next section---nor that anything
resembling a thermodynamic phase transition has been located: the
transition language used here is the minimal-integer-dimensional geometric
transition, not an order-parameter limit.
